\subsubsubsection{Case I: Logical Constraint Solving}
\label{text_cases}

\begin{caseprompt}
\textbf{The Protocol of the Corrupted AI.}
A cybersecurity analyst must override a rogue AI system by entering five unique Access Keys into five slots. Each Access Key contains exactly one letter from \(\{A,B,C,D,E\}\) and one digit from \(\{1,2,3,4,5\}\). No letter or digit may repeat. Five intercepted nodes describe the sequence, but exactly one node is corrupted: every statement in the corrupted node is false, while every statement in each intact node is true.

\begin{description}[leftmargin=1.8em,itemsep=0.2em,topsep=0.3em]
  \item[Node 1.] \(A\) is immediately to the left of the slot containing digit \(5\). Digit \(3\) is located in Slot 5.
  \item[Node 2.] \(C\) is paired with an even digit. The slot containing \(E\) is exactly two slots to the right of the slot containing digit \(4\).
  \item[Node 3.] Digit \(2\) is adjacent to the slot containing \(D\). \(B\) is neither in Slot 1 nor Slot 5.
  \item[Node 4.] The sum of the digits in Slot 1 and Slot 2 equals the digit in Slot 3. The digit paired with \(A\) is strictly less than the digit paired with \(B\).
  \item[Node 5.] \(D\) is somewhere to the left of \(C\). Digit \(1\) is paired directly with \(E\). Slot 4 contains an odd digit.
\end{description}

\textbf{Mission.} Identify the corrupted node and determine the exact letter-digit sequence for all five slots.
\end{caseprompt}

\begin{caseresponse}
\textbf{Corrupted node.} Node 5 is the corrupted node. All statements in Node 5 are false, while all statements in Nodes 1--4 are true.

\textbf{Exact sequence.}
\begin{center}
\begin{tabular}{ccc}
\toprule
\textbf{Slot} & \textbf{Letter} & \textbf{Digit} \\
\midrule
1 & C & 4 \\
2 & A & 1 \\
3 & E & 5 \\
4 & B & 2 \\
5 & D & 3 \\
\bottomrule
\end{tabular}
\end{center}

\textbf{Verification.}
\begin{center}
\begin{tabular}{p{0.13\linewidth}p{0.77\linewidth}}
\toprule
\textbf{Node} & \textbf{Status and evidence} \\
\midrule
Node 1 & True. \(A\) in Slot 2 is immediately to the left of digit \(5\) in Slot 3; digit \(3\) is in Slot 5. \\
Node 2 & True. \(C\) is paired with digit \(4\), which is even; \(E\) in Slot 3 is exactly two slots to the right of digit \(4\) in Slot 1. \\
Node 3 & True. Digit \(2\) in Slot 4 is adjacent to \(D\) in Slot 5; \(B\) is in Slot 4, so it is neither in Slot 1 nor Slot 5. \\
Node 4 & True. The Slot 1 and Slot 2 digits sum to the Slot 3 digit: \(4+1=5\). The digit paired with \(A\) is \(1\), which is less than the digit paired with \(B\), \(2\). \\
Node 5 & False. \(D\) is in Slot 5 and \(C\) is in Slot 1, so \(D\) is not left of \(C\). Digit \(1\) is paired with \(A\), not \(E\). Slot 4 contains digit \(2\), which is even rather than odd. \\
\bottomrule
\end{tabular}
\end{center}

The assignment satisfies the uniqueness constraint for all letters and digits and leaves exactly one corrupted node.
\end{caseresponse}

\captionof{figure}{\textbf{Text case for logical constraint solving.} Keye-VL-2.0 solves a multi-constraint symbolic puzzle by identifying the single corrupted evidence node, constructing a valid slot assignment, and verifying every statement against the final sequence.}
\label{fig:corrupted_ai_protocol_case}

\clearpage
